\documentclass[10pt, aspectratio=169]{beamer}
% Preámbulo modular para presentaciones Beamer (Modelación Estadística)
% Diseño y paleta institucional del Tecnológico de Monterrey

\documentclass[aspectratio=169,xcolor={dvipsnames,table}]{beamer}

\usepackage[utf8]{inputenc}
\usepackage[spanish,mexico]{babel}
\usepackage{amsmath,amssymb,amsfonts,mathtools}
\usepackage{graphicx}
\usepackage{listings}
\usepackage{booktabs}
\usepackage{tikz}
\usetikzlibrary{matrix,arrows,decorations.pathmorphing,shapes.geometric,positioning}

% Paleta Institucional Tecnológico de Monterrey
\definecolor{TecAzul}{HTML}{0039A6}       % Azul TEC: color primario institucional
\definecolor{TecAzulOscuro}{HTML}{1A2E51} % Azul marino oscuro
\definecolor{TecRojo}{HTML}{EC2661}       % Rojo de acento (paleta miTec)
\definecolor{TecGris}{HTML}{646464}       % Gris neutro para comentarios y subtítulos
\definecolor{TecFondoBloque}{HTML}{F4F6F9}% Fondo suave para bloques

% Configuración de Tema Beamer
\usetheme{Boadilla}
\usecolortheme[named=TecAzulOscuro]{structure}
\setbeamercolor{alerted text}{fg=TecRojo}
\setbeamercolor{palette primary}{bg=TecAzulOscuro,fg=white}
\setbeamercolor{palette secondary}{bg=TecAzul,fg=white}
\setbeamercolor{palette tertiary}{bg=TecAzulOscuro!80!black,fg=white}
\setbeamercolor{title}{fg=TecAzulOscuro,bg=TecFondoBloque}
\setbeamercolor{frametitle}{fg=TecAzulOscuro,bg=TecFondoBloque}

% Bloques personalizados elegantes
\setbeamercolor{block title}{bg=TecAzulOscuro,fg=white}
\setbeamercolor{block body}{bg=TecFondoBloque,fg=black}
\setbeamercolor{block title alerted}{bg=TecRojo,fg=white}
\setbeamercolor{block body alerted}{bg=TecRojo!8,fg=black}
\setbeamercolor{block title example}{bg=TecAzul,fg=white}
\setbeamercolor{block body example}{bg=TecAzul!8,fg=black}

% Redondear bordes y añadir sombra sutil
\setbeamertemplate{blocks}[rounded][shadow=true]
\setbeamertemplate{navigation symbols}{}

% Pie de página limpio con número de diapositiva
\setbeamertemplate{footline}{
  \leavevmode%
  \hbox{%
  \begin{beamercolorbox}[wd=.7\paperwidth,ht=2.25ex,dp=1ex,left]{author in head/foot}%
    \hspace*{2ex}\insertshortauthor~~(\insertshortinstitute) \hfill \insertshorttitle
  \end{beamercolorbox}%
  \begin{beamercolorbox}[wd=.3\paperwidth,ht=2.25ex,dp=1ex,right]{date in head/foot}%
    \insertshortdate{}\hspace*{2em}
    \insertframenumber{} / \inserttotalframenumber\hspace*{2ex}
  \end{beamercolorbox}}%
  \vskip0pt%
}

% Configuración de Listings de Python para visualización pedagógica de código
\lstset{
    language=Python,
    basicstyle=\ttfamily\footnotesize,
    keywordstyle=\color{TecAzulOscuro}\bfseries,
    stringstyle=\color{TecRojo},
    commentstyle=\color{TecGris}\itshape,
    showstringspaces=false,
    breaklines=true,
    frame=lines,
    rulecolor=\color{TecAzulOscuro!30},
    backgroundcolor=\color{TecFondoBloque},
    aboveskip=0.5em,
    belowskip=0.5em,
    literate={á}{{\'a}}1 {é}{{\'e}}1 {í}{{\'i}}1 {ó}{{\'o}}1 {ú}{{\'u}}1
             {Á}{{\'A}}1 {É}{{\'E}}1 {Í}{{\'I}}1 {Ó}{{\'O}}1 {Ú}{{\'U}}1
             {ñ}{{\~n}}1 {Ñ}{{\~N}}1 {¿}{{?`}}1 {¡}{{!`}}1
}

% Metadatos institucionales por defecto
\title[Modelación Estadística]{Modelación Estadística para Ciencia de Datos e Ingeniería}
\author[J. Castillo Colmenares]{Juliho David Castillo Colmenares}
\institute[Tec de Monterrey]{Tecnológico de Monterrey \\ Escuela de Ingeniería y Ciencias}
\date{\today}

% Comandos y macros estadísticas compartidas para las presentaciones Beamer (Metropolis)

% Conjuntos numéricos básicos
\newcommand{\R}{\mathbb{R}}
\newcommand{\N}{\mathbb{N}}
\newcommand{\Q}{\mathbb{Q}}
\newcommand{\Z}{\mathbb{Z}}
\newcommand{\set}[1]{\left\{ #1 \right\}}
\newcommand{\abs}[1]{\left|#1\right|}
\newcommand{\norm}[1]{\left\| #1 \right\|}

% Comandos estadísticos y probabilísticos del curso
\newcommand{\Var}{\operatorname{Var}}
\newcommand{\cov}{\operatorname{Cov}}
\newcommand{\corr}{\rho}
\newcommand{\comb}[2]{\begin{pmatrix} #1 \\ #2 \end{pmatrix}}
\newcommand{\s}{\sigma}
\newcommand{\particion}{\mathfrak{P}}
\newcommand{\card}[1]{\#\left( #1 \right)}
\newcommand{\seq}[3]{\set{#1}_{#2}^{#3}}
\newcommand{\E}{\mathbb{E}}
\newcommand{\Prob}{\mathbb{P}}

% Entornos pedagógicos y cajas en español académico natural
\newenvironment{discusion}{%
  \begin{alertblock}{Pregunta de Discusión}%
}{%
  \end{alertblock}%
}

\newenvironment{reflexion}{%
  \begin{alertblock}{Reflexión para el Aula}%
}{%
  \end{alertblock}%
}

\newenvironment{paradoja}{%
  \begin{alertblock}{Paradoja de Exploración}%
}{%
  \end{alertblock}%
}

\newenvironment{motivacion}{%
  \begin{exampleblock}{Motivación: Ciencia de Datos en la Práctica}%
}{%
  \end{exampleblock}%
}

\newenvironment{retoclase}{%
  \begin{block}{Reto Rápido en Equipo}%
}{%
  \end{block}%
}


\title[PDF y Soporte Continuo]{Sección 04.01: PDF y Soporte Continuo}
\subtitle{Función de Densidad, Esperanza Continua y LOTUS}
\author[J. Castillo Colmenares]{Juliho Castillo Colmenares}
\institute{Tecnológico de Monterrey}
\date{\vspace{-1.2cm}}

\begin{document}

% Slide 1: Portada
\begin{frame}[plain]
  \titlepage
\end{frame}

% Slide 2: Hoja de Ruta
\begin{frame}{Hoja de Ruta --- Capítulo 04: Variables Aleatorias Continuas}
  \footnotesize
  ¿Dónde nos encontramos en el desarrollo modular del capítulo? \pause
  \begin{itemize}\itemsep=0.08em
    \item \textbf{04.01} Función de Densidad (PDF) y Soporte Continuo \emph{(Hoy)} \pause
    \item \textbf{04.02} Función de Distribución Acumulada (CDF) Continua y Cuantiles \pause
    \item \textbf{04.03} Esperanza Matemática, Varianza y LOTUS Continuo \pause
    \item \textbf{04.04} Distribución Uniforme Continua \pause
    \item \textbf{04.05} Distribución Normal y Puntaje $Z$ \pause
    \item \textbf{04.06} Distribución Exponencial y Procesos sin Memoria \pause
    \item \textbf{04.07} Distribuciones Gamma, Beta y Weibull
  \end{itemize}
  \vspace{-0.05cm}
  \begin{block}{Objetivo de la Sesión}
    Introducir la función de densidad de probabilidad (PDF) como análogo continuo de la PMF, dominar la condición de normalización $\int f = 1$, calcular probabilidades mediante integración, y establecer los fundamentos de la esperanza continua y el Teorema LOTUS.
  \end{block}
\end{frame}

% Slide 3: Motivación
\begin{frame}{De Discreto a Continuo: La Función de Densidad}
  \footnotesize
  \begin{motivacion}[¿Por qué variables continuas?]
    En el Capítulo 03 trabajamos con variables que toman valores discretos: número de aciertos, defectos, llamadas. Pero muchos fenómenos varían continuamente: tiempo de espera, errores de medición, voltaje en un circuito, peso de un producto. Para estas variables, la masa de probabilidad se ``distribuye'' sobre un continuo, y la \emph{densidad de probabilidad} $f(x)$ codifica cómo se concentra esa masa.
  \end{motivacion}

  \vspace{0.15cm}
  \pause
  \begin{itemize}\itemsep=0.1cm
    \item \textbf{Continuidad:} $X$ toma valores en un intervalo (o unión de intervalos), no en un conjunto discreto. \pause
    \item \textbf{Probabilidad como área:} $P(a \le X \le b) = \int_{a}^{b} f(x)\,dx$, no como suma de masas puntuales. \pause
    \item \textbf{Condición técnica:} $\int_{-\infty}^{\infty} f(x)\,dx = 1$ garantiza que la masa total es la unidad.
  \end{itemize}
\end{frame}

% Slide 4: Definición Formal de PDF
\begin{frame}{Definición Formal de la Función de Densidad de Probabilidad}
  \footnotesize
  Una variable aleatoria $X$ es \emph{continua} si existe una función integrable $f: \mathbb{R} \to [0, \infty)$ tal que para todo intervalo $[a, b]$: \pause

  \vspace{0.1cm}
  \begin{block}{Función de Densidad de Probabilidad (PDF)}
    $$P(a \le X \le b) = \int_{a}^{b} f(x)\,dx$$
    con las dos condiciones axiomáticas: $f(x) \ge 0$ para todo $x$ y $\int_{-\infty}^{\infty} f(x)\,dx = 1$.
  \end{block}

  \vspace{0.1cm}
  \pause
  \begin{alertblock}{Propiedad Fundamental: $P(X = x) = 0$}
    \begin{itemize}\itemsep=0.05cm
      \item Para una variable continua, cualquier valor puntual tiene probabilidad nula.
      \item Las probabilidades se calculan sobre \emph{intervalos}, no sobre puntos.
      \item Por lo tanto, $P(a \le X \le b) = P(a < X \le b) = P(a \le X < b) = P(a < X < b)$.
    \end{itemize}
  \end{alertblock}
\end{frame}

% Slide 5: Soporte Continuo
\begin{frame}{Soporte Continuo y Complementos del Soporte}
  \footnotesize
  El \emph{soporte} de una variable continua es el conjunto de puntos donde la PDF es estrictamente positiva: \pause

  \vspace{0.1cm}
  \begin{block}{Soporte Continuo $\mathcal{S}_X$}
    $$\mathcal{S}_X = \{x \in \mathbb{R} : f(x) > 0\}$$
    Para variables continuas, el soporte puede ser un intervalo acotado $[a, b]$, semi-acotado $[a, \infty)$ o $(-\infty, b]$, o toda la recta real $\mathbb{R}$.
  \end{block}

  \vspace{0.1cm}
  \pause
  \begin{alertblock}{Ejemplos de Soporte}
    \begin{itemize}\itemsep=0.05cm
      \item Uniforme $U(0, 1)$: soporte $[0, 1]$ (acotado, cerrado).
      \item Exponencial $\text{Exp}(\lambda)$: soporte $[0, \infty)$ (semi-acotado, cola derecha).
      \item Normal $N(\mu, \sigma^{2})$: soporte $\mathbb{R}$ (no acotado, dos colas).
    \end{itemize}
  \end{alertblock}
\end{frame}

% Slide 6: CDF Continua
\begin{frame}{Función de Distribución Acumulada Continua (CDF)}
  \footnotesize
  La CDF se obtiene integrando la PDF desde $-\infty$ hasta $x$: \pause

  \vspace{0.1cm}
  \begin{block}{CDF Continua}
    $$F(x) = P(X \le x) = \int_{-\infty}^{x} f(t)\,dt$$
    Por el Teorema Fundamental del Cálculo, $f(x) = F'(x)$ siempre que $F$ sea diferenciable. La CDF es monótona creciente con $F(-\infty) = 0$ y $F(+\infty) = 1$.
  \end{block}

  \vspace{0.1cm}
  \pause
  \begin{alertblock}{Cálculo de Probabilidades por Intervalos}
    \begin{itemize}\itemsep=0.05cm
      \item $P(a \le X \le b) = F(b) - F(a) = \int_{a}^{b} f(t)\,dt$.
      \item $P(X > a) = 1 - F(a)$.
      \item $P(X \le b) = F(b)$ por definición.
    \end{itemize}
  \end{alertblock}
\end{frame}

% Slide 7: Esperanza Continua
\begin{frame}{Esperanza Matemática Continua: $\E[X] = \int x f(x)\,dx$}
  \footnotesize
  La esperanza de una variable continua se calcula integrando $x \cdot f(x)$ sobre todo el soporte: \pause

  \vspace{0.1cm}
  \begin{block}{Definición de Esperanza Continua}
    $$\E[X] = \int_{-\infty}^{\infty} x \cdot f(x)\,dx$$
    siempre que la integral exista (condición de integrabilidad absoluta: $\int |x| f(x)\,dx < \infty$).
  \end{block}

  \vspace{0.1cm}
  \pause
  \begin{alertblock}{Varianza Continua}
    \begin{itemize}\itemsep=0.05cm
      \item $\Var(X) = \E[(X - \mu)^{2}] = \int (x - \mu)^{2} f(x)\,dx$.
      \item Forma computacional: $\Var(X) = \E[X^{2}] - (\E[X])^{2}$ donde $\E[X^{2}] = \int x^{2} f(x)\,dx$.
      \item Desviación estándar: $\sigma = \sqrt{\Var(X)}$.
    \end{itemize}
  \end{alertblock}
\end{frame}

% Slide 8: LOTUS Continuo
\begin{frame}{Teorema LOTUS Continuo: $\E[g(X)] = \int g(x) f(x)\,dx$}
  \footnotesize
  Para cualquier función $g: \mathbb{R} \to \mathbb{R}$, el Teorema LOTUS (Law of the Unconscious Statistician) permite calcular la esperanza de $g(X)$ sin necesidad de encontrar la distribución de $g(X)$: \pause

  \vspace{0.1cm}
  \begin{block}{LOTUS Continuo}
    $$\E[g(X)] = \int_{-\infty}^{\infty} g(x) \cdot f(x)\,dx$$
    bajo condiciones de regularidad (e.g., $g$ medible y $g(X)$ integrable). Caso particular: $g(x) = x$ recupera $\E[X]$; $g(x) = (x - \mu)^{2}$ recupera la varianza.
  \end{block}

  \vspace{0.1cm}
  \pause
  \begin{alertblock}{Aplicaciones Prácticas}
    \begin{itemize}\itemsep=0.05cm
      \item Transformaciones de variables: $\E[\log X]$, $\E[X^{2}]$, $\E[\sqrt{X}]$.
      \item Esperanza de funciones de pago en finanzas: $\E[\max(X - K, 0)]$ para opciones.
      \item Esperanza de indicadores continuos: $\E[\mathbb{1}_{X > c}] = P(X > c)$.
    \end{itemize}
  \end{alertblock}
\end{frame}

% Slide 9: Aplicaciones
\begin{frame}{Aplicaciones en Ingeniería, Ciencias y Ciencia de Datos}
  \footnotesize
  Las variables continuas y sus PDFs son omnipresentes en modelado moderno: \pause

  \vspace{0.15cm}
  \begin{itemize}\itemsep=0.12cm
    \item \textbf{Telecomunicaciones:} Tiempo de espera de paquetes, ruido gaussiano en canales, y modelos de interferencia. \pause
    \item \textbf{Confiabilidad Industrial:} Tiempo hasta la falla de componentes mecánicos y electrónicos, modelado exponencial/Weibull. \pause
    \item \textbf{Física Experimental:} Errores de medición con distribución Normal, propagación de incertidumbres. \pause
    \item \textbf{Finanzas y Riesgo:} Movimientos Brownianos para precios de activos, modelos de valuación de opciones.
  \end{itemize}
\end{frame}

% Slide 12A-1: Lab Python (Bloque 1)
\begin{frame}[fragile]{Laboratorio en Python: PDF y Normalización}
  \scriptsize
  Verificación computacional de la condición $\int f = 1$ para Exponential, Rayleigh y Cuadrática (\texttt{04.01\_pdf\_and\_support.py}):
  \vspace{0.05cm}
  {\lstset{basicstyle=\fontsize{5pt}{6pt}\selectfont\ttfamily}
  \lstinputlisting[firstline=21, lastline=44]{../../code/04_variables_aleatorias_continuas/04.01_pdf_and_support.py}}
\end{frame}

% Slide 12A-2: Lab Python (Bloque 2)
\begin{frame}[fragile]{Laboratorio en Python: CDF y Probabilidades de Intervalo}
  \scriptsize
  Cálculo numérico de CDF y probabilidades $P(a \le X \le b)$ con `scipy.integrate.quad`:
  \vspace{0.02cm}
  {\lstset{basicstyle=\fontsize{4.5pt}{5.5pt}\selectfont\ttfamily}
  \lstinputlisting[firstline=47, lastline=76]{../../code/04_variables_aleatorias_continuas/04.01_pdf_and_support.py}}
\end{frame}

% Slide 12B: Lab Python (Bloque 3)
\begin{frame}[fragile]{Laboratorio en Python: LOTUS Continuo y Momentos}
  \scriptsize
  Verificación de LOTUS: $\E[g(X)] = \int g(x) f(x)\,dx$ para Exponential y Rayleigh:
  \vspace{0.02cm}
  {\lstset{basicstyle=\fontsize{4.5pt}{5.5pt}\selectfont\ttfamily}
  \lstinputlisting[firstline=85, lastline=112]{../../code/04_variables_aleatorias_continuas/04.01_pdf_and_support.py}}
\end{frame}

% Slide 12C: Salida Terminal
\begin{frame}[fragile]{Laboratorio en Python: Salida en Terminal}
  \scriptsize
  Salida estándar generada al ejecutar el script del laboratorio en Python:
  \vspace{0.02cm}
  \begin{block}{Salida Estándar en Consola}
    \fontsize{4pt}{5pt}\selectfont
    \begin{verbatim}
=== Block 1: PDF Validation & Normalization ===
  Exponential(0.5):         integral = 1.00000000
  Rayleigh(1.0):            integral = 1.00000000
  Quadratic(3/10):          integral = 1.00000000
  SciPy Rayleigh match: True

=== Block 2: CDF & Interval Probabilities ===
Exponential(0.5) CDF:
  F(0.5)=0.221199 | F(1.0)=0.393469 | F(2.0)=0.632121
P(0<=X<=1)=0.393469 | P(1<=X<=3)=0.383400
Rayleigh CDF: F(0.5)=0.117503 | F(1.0)=0.393469
Quadratic: F(-0.5)=0.2125 | F(0)=0.5 | F(0.5)=0.7875

=== Block 3: Continuous LOTUS & Moments ===
Exponential(0.5): E[X]=2.0 | E[X^2]=8.0 | Var=4.0
Rayleigh(1.0): E[X]=1.253314 | Var=0.429204
LOTUS E[sqrt(X)]: numerical=1.253314 | SciPy=1.253314
    \end{verbatim}
  \end{block}
\end{frame}

% Slide 13A: Ejercicio Nivel 1 (Enunciado)

% Slide 13B: Ejercicio Nivel 1 (Resolución)

% Slide 14A: Ejercicio Nivel 2 (Enunciado)

% Slide 14B: Ejercicio Nivel 2 (Resolución)

% Slide 15A: Ejercicio Nivel 3 (Enunciado)

% Slide 17: Cierre
\begin{frame}{Cierre de Sección y Síntesis sobre Variables Continuas}
  \begin{reflexion}[El Puente de lo Discreto a lo Continuo]
    La función de densidad de probabilidad $f(x)$ es el análogo continuo de la PMF: en lugar de masa en puntos, distribuye probabilidad sobre un continuo. Las propiedades formales son análogas (no-negatividad, integral 1), pero las técnicas de cálculo requieren integración en lugar de sumas.
  \end{reflexion}

  \vspace{0.2cm}
  \pause
  \begin{block}{Lecciones Clave de la PDF}
    \begin{itemize}\itemsep=0.08cm
      \item \textbf{Normalización:} $\int f = 1$ garantiza la masa total; las PDFs son relativas, no absolutas.
      \item \textbf{Soporte:} Determina dónde concentrar el cálculo; puede ser acotado, semi-acotado, o no acotado.
      \item \textbf{LOTUS:} $\E[g(X)] = \int g(x) f(x)\,dx$ evita derivar la distribución de $g(X)$.
    \end{itemize}
  \end{block}
\end{frame}

% Slide 18: Perspectiva Modular
\begin{frame}{Perspectiva Modular --- El Próximo Reto}
  \scriptsize
  \begin{retoclase}[Para la próxima sesión: CDF Continua y Cuantiles]
    Acabamos de introducir la PDF como objeto fundamental de las variables continuas. La CDF continua, definida como $F(x) = \int_{-\infty}^{x} f(t)\,dt$, es la herramienta complementaria que permite calcular probabilidades acumuladas y definir cuantiles. También estudiaremos la inversión de la CDF, fundamental para la generación de números aleatorios y los métodos de Monte Carlo.
  \end{retoclase}

  \vspace{0.08cm}
  \pause
  \begin{alertblock}{Preparación Recomendada}
    \begin{itemize}\itemsep=0.06cm
      \item Repasa el Teorema Fundamental del Cálculo y la integración por partes.
      \item Reflexiona sobre cómo encontrar $x$ tal que $F(x) = p$ (cuantil) y su utilidad en inferencia.
    \end{itemize}
  \end{alertblock}
\end{frame}

\end{document}
